\section*{Capítulo I, Problema XIV}

\textbf{Problema:}

14. Demuestra que si $\theta$ es el ángulo entre $A$ y $B$, entonces:
\[
    \|A - B\|^2 = \|A\|^2 + \|B\|^2 - 2\|A\|\|B\| \cos \theta.
\]

\textbf{Solución:}

Por definición de norma, tenemos:
\begin{align*}
    \|A - B\|^2 &= (A - B) \cdot (A - B) \\
    &= A \cdot A - 2(A \cdot B) + B \cdot B \\
    &= \|A\|^2 - 2(A \cdot B) + \|B\|^2.
\end{align*}

Sabemos que el producto escalar $A \cdot B$ se puede expresar en términos de las normas de $A$ y $B$ y el coseno del ángulo entre ellos:
\[
    A \cdot B = \|A\|\|B\| \cos \theta.
\]

Sustituyendo esta expresión en la ecuación anterior:
\begin{align*}
    \|A - B\|^2 &= \|A\|^2 - 2(\|A\|\|B\| \cos \theta) + \|B\|^2 \\
    &= \|A\|^2 + \|B\|^2 - 2\|A\|\|B\| \cos \theta.
\end{align*}

Por lo tanto, se cumple la relación dada.